\section*{अध्याय \ref{ch:b2:unicellular-diversity} --- एककोशिकीय जीवों की विविधता}

\begin{solution}{exo:b2:unicellular-diversity:1}
\omterm{def:b2:unicellular-diversity:unicellular}{एककोशिकीय}: एक ही कोशिका सारे कार्य निभाती है और अकेली जीती है;
\omterm{def:b2:unicellular-diversity:unicellular}{निवही}: एक जैसी कोशिकाएँ जुड़ी रहती हैं पर हर एक अकेली जी सकती है;
\omterm{def:b2:unicellular-diversity:unicellular}{बहुकोशिकीय}: कई प्रकार की कोशिकाएँ, एक-दूसरे पर निर्भर, और जनन केवल
जनन वंशक्रम करता है। \emph{E.~coli} और खमीर \omterm{def:b2:unicellular-diversity:unicellular}{एककोशिकीय} हैं (मुकुल
अलग होकर अकेला जीता है); \emph{Gonium} \omterm{def:b2:unicellular-diversity:unicellular}{निवही} है (16 एक जैसी
कोशिकाएँ, हर एक नया निवह बसाने में समर्थ); \emph{Volvox} \omterm{def:b2:unicellular-diversity:unicellular}{बहुकोशिकीय} है
(नश्वर कायिक कोशिकाएँ, कुछ जनन कोशिकाएँ)।
\end{solution}

\begin{solution}{exo:b2:unicellular-diversity:2}
$S/V = 3/r$: गोलाणु के लिए \qty{6}{\micro m^{-1}},
\omterm{def:b2:unicellular-diversity:protist}{प्रोटिस्ट} के लिए \qty{0.06}{\micro m^{-1}}, यानी सौ गुना अंतर।
प्रति इकाई आयतन गोलाणु तेज़ी से श्वसन करता है: उसके कोशिकाद्रव्य का हर
भाग उस पृष्ठ से माइक्रोमीटर के एक अंश भर दूर है जिससे ऑक्सीजन और भोजन
भीतर आते हैं, जबकि \omterm{def:b2:unicellular-diversity:protist}{प्रोटिस्ट} का भीतरी भाग प्रति इकाई आयतन सौ गुना कम
पृष्ठ से पलता है।
\end{solution}

\begin{solution}{exo:b2:unicellular-diversity:3}
केंद्रक नहीं (एक \omterm{def:b2:unicellular-diversity:prokaryote}{केंद्रकाभ}) बनाम केंद्रक आवरण; न सूत्रकणिकाएँ न
अंतःझिल्लियाँ, श्वसन शृंखला प्लाज़्मा झिल्ली में बैठी हुई, बनाम
कोशिकांग; ऐसा कशाभ जो फ्लैजेलिन की घूमती दृढ़ कुंडली है और प्रोटॉनों से
घुमाया जाता है, बनाम $9+2$ पक्ष्माभ जो डाइनीन और ATP से मुड़ता है; और
आकार भी (\qty{1}{\micro m} बनाम कई दहाई) तथा पेप्टाइडोग्लाइकैन भित्ति।
कोई जीवाणु भक्षण नहीं कर सकता: उसकी भित्ति किसी कण को निगलने से रोकती
है, जो कोई अमीबा अपने कूटपादों से कर लेता है।
\end{solution}

\begin{solution}{exo:b2:unicellular-diversity:4}
पक्ष्माभ: गमन और भोजन को मुख की ओर बुहारना; मुखीय नाली: कणों को गुहिका
तक पहुँचाना; भोजन-रसधानी: निगले हुए जीवाणुओं को लयनकाय के एंजाइमों से
पचाना; \omterm{prop:b2:unicellular-diversity:osmoregulation}{संकुचनशील रसधानी}: परासरण से भीतर आते जल को बाहर फेंकना;
गुरुकेंद्रक: कामकाजी केंद्रक, जिसमें रोज़मर्रा के जीवन के लिए जीन की
सैकड़ों प्रतियाँ अनुलेखित होती हैं; लघुकेंद्रक: जननिक द्विगुणित केंद्रक,
जो अर्धसूत्री विभाजन और संयुग्मन में काम आता है।
\end{solution}

\begin{solution}{exo:b2:unicellular-diversity:5}
$t = x^2/2D$: जीवाणु के लिए $(10^{-6})^2/10^{-9} = \qty{1}{ms}$;
अमीबा के लिए $(5\times 10^{-4})^2/10^{-9} = \qty{250}{s}$, यानी लगभग
चार मिनट। अमीबा उपापचयजों को बाँटने के लिए विसरण पर नहीं टिक सकता: उसे
कोशिकाद्रव्य का प्रवाह और कोशिका-कंकाल के साथ मोटर-चालित अभिगमन चाहिए,
जिनके बिना जीवाणु काम चला लेता है।
\end{solution}

\begin{solution}{exo:b2:unicellular-diversity:6}
$Re = \rho vL/\eta$: \emph{Paramecium} $10^3\times 10^{-3}\times
2\times 10^{-4}/10^{-3} = 0.2$; जीवाणु $10^3\times 2\times
10^{-5}\times 2\times 10^{-6}/10^{-3} = 4\times 10^{-5}$; तैराक
$10^3\times 1\times 2/10^{-3} = 2\times 10^6$। $Re \gg 1$ पर जड़त्व
तैराक को आघात के बाद आगे ले जाता है; $Re \ll 1$ पर श्यान कर्षण
जीवाणु को एक नैनोमीटर के भीतर रोक देता है, इसलिए वह तभी चलता है जब
स्पंदन करता है।
\end{solution}

\begin{solution}{exo:b2:unicellular-diversity:7}
$v_s = 2r^2\Delta\rho g/9\eta = 2\times 10^{-10}\times 100\times
9.81/(9\times 10^{-3}) = \qty{2.2e-5}{m/s}$; \qty{50}{m} में
$50/2.2\times 10^{-5} = \qty{2.3e6}{s}$, यानी 27 दिन लगते हैं। त्रिज्या
आधी करने पर $v_s$ 4 से भाग जाता है: 108 दिन। $\Delta\rho$ घटाकर
\qty{20}{kg/m^3} करने पर वह 5 से भाग जाता है: 135 दिन।
\end{solution}

\begin{solution}{exo:b2:unicellular-diversity:8}
$V = \tfrac{4}{3}\pi\times 100\times 25\times 25 =
\qty{2.6e5}{\micro m^3}$। \qty{1800}{s} में: $2.6\times 10^5/1800 =
\qty{145}{\micro m^3/s}$; और चूँकि $\qty{1}{\micro m^3} =
\qty{1e-15}{L}$, यह \qty{1.5e-13}{L/s} है।
\end{solution}

\begin{solution}{exo:b2:unicellular-diversity:9}
\qty{1}{mm} में दुगुना होना प्रति माइक्रोमीटर लगभग \qty{0.07}{\%} की
वृद्धि है (क्योंकि $2^{1/1000} - 1 = 0.0007$), इसलिए कोशिका भर में
\qty{0.07}{\%}। \qty{1}{s} की दौड़ \qty{20}{\micro m} तय करती है:
\qty{1.4}{\%}। \qty{0.07}{\%} का अंतर सुलझाने के लिए कोशिका को हर सिरे
पर कोई $(1/0.0007)^2 = 2\times 10^6$ अणु गिनने पड़ते, जो एक सेकंड में
उसके कुछ हज़ार ग्राहियों से बँधने वाले अणुओं से कहीं अधिक है;
\qty{1.4}{\%} के लिए लगभग \num{5000} चाहिए, जो संभव है। कालिक तुलना
एक असंभव स्थानिक मापन को संभव बना देती है, क्योंकि समाकलन का काम दौड़
पर छोड़ दिया जाता है।
\end{solution}

\begin{solution}{exo:b2:unicellular-diversity:10}
पृष्ठ $4\pi r^2$ जिसमें $r = \qty{250}{\micro m}$: \qty{7.9e5}{\micro
m^2}; कोशिकाद्रव्यी आयतन $\approx$ पृष्ठ $\times$ \qty{2}{\micro m}
$= \qty{1.6e6}{\micro m^3}$; अनुपात \qty{0.5}{\micro m^{-1}}।
\emph{E.~coli} के लिए, एक बेलन: $S/V = 2\pi r L/\pi r^2 L = 2/r =
\qty{4}{\micro m^{-1}}$। वह दानव \emph{E.~coli} से केवल आठ गुना ही
ख़राब है, 500 गुना नहीं, क्योंकि उसके कोशिकाद्रव्य का हर बिंदु पृष्ठ से
\qty{2}{\micro m} के भीतर है। वह रसधानी नाइट्रेट सँजोए रखती है, यानी
उसका इलेक्ट्रॉन ग्राही, ताकि वह अनॉक्सी अवसाद में उन विरल मथनों के बीच
सल्फाइड पर श्वसन कर सके जो नाइट्रेट लाती हैं
(\cref{ch:b2:microbial-metabolism})।
\end{solution}

\begin{solution}{exo:b2:unicellular-diversity:11}
कोई कोशिका एक साथ तैर और विभाजित नहीं हो सकती। 16 कोशिकाओं का निवह
विभाजन के घंटों में तैरना रोक दे तो स्टोक्स के नियम से उपेक्षणीय
दूरी डूबता है (उसकी त्रिज्या कुछ दहाई माइक्रोमीटर है, उसकी डूबन-चाल
माइक्रोमीटर प्रति सेकंड); पूरा निवह विभाजित हो सकता है और कुछ नहीं
खोता। \num{10000} कोशिकाओं का \qty{500}{\micro m} त्रिज्या वाला गोला सौ
गुना तेज़ डूबता --- मिलिमीटर प्रति सेकंड, घंटे भर में मीटर --- और
प्रकाश से बाहर निकल जाता: इसलिए अधिकांश कोशिकाओं को तैरता रखना और कुछ
को विभाजित होने देना लाभ का सौदा है। जनन कोशिकाएँ इसलिए बड़ी हैं कि
उन्हें बिना खाए, क्रमिक विभाजनों से, हज़ारों कोशिकाओं का पूरा
पुत्री-गोला बनाने के संसाधन ढोने पड़ते हैं, और उन्हें काया ही रसद देती
है।
\end{solution}

\begin{solution}{exo:b2:unicellular-diversity:12}
क्लेड कोई पूर्वज और उसकी सारी संतति है; स्तर वह संगठन-स्तर है जिस तक
कई वंशक्रम पहुँचे हैं। \omterm{def:b2:unicellular-diversity:protist}{प्रोटिस्टों} में --- हरी शैवाल (स्थलीय पौधों के
साथ), डायटम (भूरी शैवालों के साथ), पक्ष्माभी (डाइनोफ्लैजेलेट और
\emph{Plasmodium} के साथ), अमीबा (प्राणियों और कवकों के निकट) --- केवल
अकेली जी जाती सुकेंद्रकी कोशिका साझी है, और उनका अंतिम साझा पूर्वज हम
समेत सारे सुकेंद्रकियों का पूर्वज है: यानी एक स्तर। ``\omterm{def:b2:unicellular-diversity:prokaryote}{प्राक्केंद्रकी}''
जीवाणुओं और आर्किया को उस बात से जोड़ता है जो उनमें नहीं है; आर्किया
जीवाणुओं से अधिक सुकेंद्रकियों के निकट हैं, इसलिए वह शब्द भी एक स्तर
है। सायनोजीवाणु उस एक पूर्वज से उतरे हैं जिसने ऑक्सीजनी प्रकाशसंश्लेषण
निकाला, और उसकी सारी संतति उनमें है (हरितलवक वंशक्रम को छोड़कर): यानी
एक क्लेड।
\end{solution}

\begin{solution}{pb:b2:unicellular-diversity:1}
\textbf{1.} $1\times 1\times 0.1 = \qty{0.1}{mm^3} = \qty{0.1}{\micro L}$।
\textbf{2.} $24/\qty{0.1}{\micro L} = 240$ प्रति \unit{\micro L} $=
\qty{2.4e5}{}$ कोशिकाएँ प्रति मिलिलिटर।
\textbf{3.} $V = \tfrac{4}{3}\pi\times 125 = \qty{524}{\micro m^3}$;
$S/V = 3/5 = \qty{0.6}{\micro m^{-1}}$।
\textbf{4.} प्रति मिलिलिटर $2.4\times 10^5\times 524 = \qty{1.26e8}{\micro m^3}$,
और $\qty{1}{mL} = \qty{1e12}{\micro m^3}$: यानी अंश
$1.3\times 10^{-4}$, \qty{0.013}{\%}।
\textbf{5.} $\qty{1000}{m^3} = \qty{1e9}{mL}$: $2.4\times 10^{14}$
कोशिकाएँ।
\textbf{6.} सायनोजीवाणु कहीं छोटी कोशिकाओं (कुछ माइक्रोमीटर) का तंतु
है, नीला-हरा और बिना किसी स्पष्ट हरितलवक या केंद्रक के, और वह कशाभों से
तैरता नहीं; शैवाल दो कशाभों, एक नेत्रबिंदु और एक प्याले-आकार के
हरितलवक वाली अकेली अंडाकार कोशिका है।
\textbf{7.} $384/24 = 16 = 2^4$: \qty{48}{h} में चार बार दुगुना,
$T = \qty{12}{h}$।
\textbf{8.} $k = \ln 2/T = 0.693/12 = \qty{0.058}{h^{-1}}$।
\textbf{9.} $t = \ln(10^8/2.4\times 10^5)/k = \ln(417)/0.058 =
6.03/0.058 = \qty{104}{h}$, यानी लगभग 4.3 दिन।
\textbf{10.} पोषक (फ़ॉस्फ़ेट, नाइट्रेट) चुक जाते हैं; कोशिकाएँ एक-दूसरे
पर छाया डालती हैं जिससे प्रति कोशिका प्रकाश घट जाता है; चरने वाले
(पक्ष्माभी, रोटिफ़र, \emph{Daphnia}) उन पर पलकर बढ़ जाते हैं; और
\qty{1e8}{} प्रति मिलिलिटर पर कोशिकाएँ आयतन का कोई पाँच प्रतिशत भर
लेतीं और घुली हुई $\mathrm{CO_2}$ चुका देतीं।
\textbf{11.} वृद्धि केवल आधे समय: माध्य दर आधी, यानी
\qty{208}{h}, 8.7 दिन।
\textbf{12.} गिनती आबादी मापती है, अलग-अलग कोशिकाएँ नहीं: यदि कोई
कोशिका हर \qty{36}{h} में 8 पुत्रियाँ छोड़े, तो आबादी \qty{12}{h} के
तीन \omterm{def:b2:unicellular-diversity:division}{पीढ़ी कालों} में 8 $= 2^3$ गुना हो जाती
है। \omterm{def:b2:unicellular-diversity:division}{पीढ़ी काल} संख्या के दुगुने होने का माध्य काल है,
विभाजन चाहे जैसे भी समूहबद्ध हों।
\textbf{13.} $Re = 10^3\times 10^{-4}\times 10^{-5}/10^{-3} = 10^{-3}$।
\textbf{14.} $m = 1050\times 5.24\times 10^{-16} = \qty{5.5e-13}{kg}$;
$6\pi\eta r = 6\pi\times 10^{-3}\times 5\times 10^{-6} =
\qty{9.4e-8}{kg/s}$; $d = 5.5\times 10^{-13}\times 10^{-4}/9.4\times
10^{-8} = \qty{5.8e-10}{m}$: यानी आधा नैनोमीटर।
\textbf{15.} $v_s = 2\times 25\times 10^{-12}\times 50\times
9.81/(9\times 10^{-3}) = \qty{2.7e-6}{m/s}$, \qty{2.7}{\micro m/s}।
\textbf{16.} \qty{1}{m} डूबने में: $1/2.7\times 10^{-6} =
\qty{3.7e5}{s}$, 4.3 दिन; \qty{100}{\micro m/s} पर ऊपर तैरने में:
\qty{1e4}{s}, 2.8 घंटे। तैरना 37 गुना जीतता है।
\textbf{17.} $c_\infty = \qty{1e-2}{mol/m^3}$: $J = 4\pi\times
1.9\times 10^{-9}\times 5\times 10^{-6}\times 10^{-2} =
\qty{1.2e-15}{mol/s} = 7.2\times 10^8$ अणु प्रति सेकंड।
\textbf{18.} \qty{100}{pg} $= \qty{1e-10}{g} = \qty{8.3e-12}{mol}$
कार्बन, \qty{12}{h} $= \qty{43200}{s}$ में दुगुना:
\qty{1.9e-16}{mol/s} $= 1.2\times 10^8$ परमाणु प्रति सेकंड। आपूर्ति
बटा माँग: $7.2/1.2 = 6$।
\textbf{19.} आपूर्ति दस गुना बढ़ती है ($\propto r$), माँग एक
हज़ार गुना ($\propto r^3$): अनुपात $6/100 = 0.06$ रह जाता है; ऐसी
कोशिका दर के केवल छह प्रतिशत पर ही बढ़ पाती, इसलिए उस आकार की
विसरण-पोषित शैवाल मथन, भीतरी अभिगमन या कहीं धीमे उपापचय के बिना असंभव
है।
\textbf{20.} $\Delta\Pi = RT\Delta c = 8.314\times 293\times 99 =
\qty{2.4e5}{Pa}$, \qty{2.4}{bar}।
\textbf{21.} क्षेत्रफल $4\pi(5\times 10^{-6})^2 = \qty{3.14e-10}{m^2}$;
अंतर्वाह $L_p A\Delta\Pi = 10^{-14}\times 3.14\times 10^{-10}\times
2.4\times 10^5 = \qty{7.6e-19}{m^3/s} = \qty{0.76}{\micro m^3/s}$।
\textbf{22.} $524/0.76 = \qty{690}{s}$, यानी लगभग ग्यारह मिनट।
\textbf{23.} $P = \Delta\Pi\times$ अभिवाह $= 2.4\times 10^5\times
7.6\times 10^{-19} = \qty{1.8e-13}{W}$।
\textbf{24.} एक ATP: $5\times 10^4/6.02\times 10^{23} =
\qty{8.3e-20}{J}$: $1.8\times 10^{-13}/8.3\times 10^{-20} = 2.2\times
10^6$ ATP प्रति सेकंड। कार्बन स्थिरीकरण: $3\times 1.2\times 10^8 =
3.6\times 10^8$ ATP प्रति सेकंड। रसधानी उसका \qty{0.6}{\%} लेती
है --- फटने के विरुद्ध सस्ता बीमा।
\textbf{25.} जल का अंतर्वाह \qty{0.76}{\micro m^3/s} (ग्यारह मिनट में
कोशिका के आयतन जितना); \omterm{prop:b2:unicellular-diversity:osmoregulation}{संकुचनशील रसधानी} के बिना कोशिका अपना आयतन
लगभग \qty{700}{s} में दुगुना कर लेती; सूखा रहना कम से कम
$2\times 10^6$ ATP प्रति सेकंड लेता है, यानी उसका प्रकाशसंश्लेषण
कार्बन पर जितना ख़र्च करता है उसका एक प्रतिशत से भी कम।
\end{solution}
